% >>> ARCHIVO GENERADO por scripts/build_figures.py. Se sobrescribe siempre (se carga desde results_macros.tex). <<<
\renewcommand{\figcaption}[1]{\ifcsname figcap@#1\endcsname\csname figcap@#1\endcsname\else\textcolor{red}{\textbf{[??caption \detokenize{#1}]}}\fi}
% \figcaptionshort{id}: short caption for the slides (caption_slides_es; falls back to the long one)
\providecommand{\figcaptionshort}[1]{\ifcsname figcapshort@#1\endcsname\csname figcapshort@#1\endcsname\else\figcaption{#1}\fi}
\expandafter\def\csname figcap@overfitting_curves\endcsname{Accuracy de train (discontinua) y de validación (continua) por época del modelo de la entrega 1 y del de esta entrega; media de 3 semillas, banda = mínimo y máximo. La de train del modelo regularizado se mide con augmentation.}
\expandafter\def\csname figcapshort@overfitting_curves\endcsname{Train (discontinua) y validación (continua), entregas 1 y 2.}
\expandafter\def\csname figcap@dropout_curve\endcsname{Curva de validación del dropout sobre el modelo de la entrega 1: accuracy de train (discontinua) y de validación (continua), media de 3 semillas.}
\expandafter\def\csname figcapshort@dropout_curve\endcsname{Accuracy de train y de validación según el dropout.}
\expandafter\def\csname figcap@l2_curve\endcsname{Curva de validación de la penalización L2 sobre el modelo de la entrega 1: accuracy de train (discontinua) y de validación (continua), media de 3 semillas.}
\expandafter\def\csname figcapshort@l2_curve\endcsname{Accuracy de train y de validación según el factor L2.}
\expandafter\def\csname figcap@batchnorm\endcsname{BatchNorm antes o después de la activación, con la learning rate de la entrega 1 o triplicada: mean recall de validación (barras, media y desviación típica de 3 semillas) y de train (trazo discontinuo).}
\expandafter\def\csname figcapshort@batchnorm\endcsname{Mean recall de validación con y sin BatchNorm (3 semillas).}
\expandafter\def\csname figcap@augmentation\endcsname{Data augmentation sobre el modelo con BatchNorm: mean recall de validación (barras, media y desviación típica de 3 semillas) y de train (trazo discontinuo).}
\expandafter\def\csname figcapshort@augmentation\endcsname{Mean recall de validación por tipo de augmentation (3 semillas).}
\expandafter\def\csname figcap@progress\endcsname{Mean recall y accuracy de validación tras cada paso de la entrega (media de 3 semillas).}
\expandafter\def\csname figcapshort@progress\endcsname{Mean recall y accuracy de validación tras cada paso.}
\expandafter\def\csname figcap@per_class_recall\endcsname{Recall de validación por clase: modelo de la entrega 1 frente al de esta entrega (media de 3 semillas en ambos).}
\expandafter\def\csname figcapshort@per_class_recall\endcsname{Recall de validación por clase: entrega 1 frente a entrega 2.}
